%==============================================================================
\section{Rendezvous}
\label{sec:rendezvous}
%==============================================================================
Mobile agents typically run on smartphones behind NAT and do not have static IP addresses.  When an agent's IP address changes (e.g., due to network roaming), all its UDP connections break.  This section describes how agents re-establish connectivity, using two complementary mechanisms: \emph{rendezvous agents} for bootstrap and \emph{friends-based rendezvous} for subsequent reconnections.

%------------------------------------------------------------------------------
\subsection{Rendezvous Agents}
\label{sec:rv-agents}
%------------------------------------------------------------------------------
A \emph{rendezvous agent} is a lightweight GLP agent that runs at a known, static, globally-routable IP address.  Its sole purpose is to coordinate hole-punching between agents that have lost connectivity.  A rendezvous agent:
\begin{itemize}
\item Has no friends list and does not participate in the social graph.
\item Accepts cold-call connections from any agent.
\item Verifies \emph{friendship proofs} (Section~\ref{sec:friendship-proof}) to confirm that the requesting agents are friends.
\item Observes each connecting agent's public IP address from the networking layer (via the \code{peer\_address/2} system predicate), serving as an address-discovery mechanism.
\item Coordinates UDP hole-punching between the two agents by relaying their addresses.
\end{itemize}

A machine that hosts a rendezvous agent is called a \emph{GLP rendezvous server}.  In the terminology of the ICE framework~\cite{rfc8445}, the GLP rendezvous server combines the roles of STUN server (address discovery via \code{peer\_address/2}) and application-specific signaling server (candidate exchange and hole-punch coordination).  It does not serve as a TURN relay; relaying, if needed, is handled at the application level via the social graph (Section~\ref{sec:rv-fallback}).

Any individual or organization can run a GLP rendezvous server.  The architecture is federated: agents may use multiple rendezvous servers for redundancy, and the system operates without any rendezvous server when agents are BLE-adjacent or share a connected friend.

While connected, agents share with their friends which rendezvous servers they use (and any preferences or priorities) via ordinary GLP messages.  Each agent thus knows, for each friend, which rendezvous servers that friend uses.  The purpose of this exchange is convergence: the rendezvous protocol is inherently asymmetric --- when agent~$A$'s IP changes, both~$A$ and the disconnected friend~$B$ must contact the \emph{same} rendezvous agent for the matching to succeed (Section~\ref{sec:bootstrap}).  Since~$B$ cannot reach~$A$ directly, $B$ must know which rendezvous agents~$A$ uses, and the only time to learn this is while they are connected.
% \dan{What purpose does the explicit exchange of RV servers serve? if not for this, upon reconnection the user could iteratively try all its RV servers until its entire friend graph is reconnected instead of relying on a specific one. This would be more resilient too, as users could be connected over multiple RVs.}

% \udi{See clarification added above.  Both sides must contact the same RV agent for matching to succeed; $B$ learns $A$'s RV agents while connected so it knows where to send its \textit{available} message when $A$ disappears.}

% \dan{I understand that there is no way for this to work with the current protocol: this is why I suggest to change the protocol as well. The current protocol is "noisy", necessitating the exchange of all RV servers used between all peers, because peers have to be connected to the same RV for reconnection to succeed. Instead, I suggest two things: \begin{itemize}
% \item Instead of the disconnected side actively trying to reconnect to the other by querying its communicated RVs, one can instead query its own RVs -- there should be an overlap, and we can save on unnecessary protocol messages
% \item Instead of the proactive push model -- one side trying to reconnect to another aggressively -- the peer that changed its IP address could just query the RVs it has to reconnect to its friends. This is a much more simple and efficient model.
% \end{itemize}}

% \udi{Both suggestions require the RV agent to reach~$B$ without~$B$'s active participation.  But~$B$ is behind NAT --- no server can reach~$B$ unless~$B$ has first opened a NAT mapping by sending a packet to that server.  This is the fundamental constraint of NAT traversal.  The alternative requires either (a)~all agents maintain persistent connections to all RVs their friends might use (far more costly than the current design), or (b)~RV-to-RV gossip (adds inter-RV coordination, contradicting their stateless, independent design).  Our protocol minimizes~$B$'s burden: $B$ acts only when it detects~$A$ went silent, and contacts only~$A$'s known RVs.}

% \dan{I am definitely pro inter-RV coordination: that's what I am trying to convey in all of the messages. Otherwise we have no way to know if a peer is really offline (e.g. phone is off), or just temporarily offline. Inter-RV coordination would not require much, as the RVs already keep who is connected through them. It just involved adding the functionality of saving all the RVs: this could be done initally as hardcoding the list, and in the future in the initial exchange of connecting to an RV, the peer could tell the RV all the RVs it's connected to. If an RV wishes to know if a peer is connected to an RV, it just queries all RVs whether a peer is connected to them, by looking up their public key.}

% \udi{Inter-RV coordination has two problems.  First, it transforms RV agents from stateless, independent servers into a coordinated network with shared state: each RV would need to maintain a registry of connected peers, respond to queries from other RVs, know the addresses of all other RVs, and handle RV failures --- a substantial increase in complexity.  Second, and more importantly, it creates a privacy concern: in the current design, an RV agent learns that~$A$ and~$B$ are friends only at the moment they request reconnection, and retains no persistent record.  Inter-RV gossip would require each RV to maintain a registry of all connected peers and respond to queries from any other RV about whether a given public key is connected --- enabling any RV operator to map the connectivity of any agent across the entire RV network.  The current design keeps RV agents minimal and privacy-preserving.  The cost of sharing RV addresses between friends is a single GLP message per friendship --- negligible compared to the infrastructure and privacy cost of inter-RV coordination.  In short, this is a NO GO, and unless I am missing something, I suggest we end this discussion with that.}
    
    
\paragraph{Friendship proof.}
\label{sec:friendship-proof}
When two agents establish friendship, each signs a \emph{friendship attestation}: a statement binding the two agents' public keys, signed by the attesting agent's Ed25519 key.  An agent can present this attestation to a rendezvous agent to prove that it is a friend of the other agent.  The rendezvous agent verifies the signature without needing to know anything about the social graph.

%------------------------------------------------------------------------------
\subsection{Friends-Based Rendezvous}
\label{sec:friends-based-rv}
%------------------------------------------------------------------------------
A \emph{friends-based rendezvous} occurs when a common friend that is currently connected to both sides coordinates the hole-punch via ordinary GLP messages.  Unlike a rendezvous agent, a friend acting in this role does not have a static IP address (smartphone agents never do), and is not dedicated infrastructure---it is simply a friend that happens to be connected to both parties.

Through the social graph, each agent knows the friends of each of its friends.  When agents~$A$ and~$B$ lose connectivity, each can identify common friends that may still be connected to both sides and request their help.

Friends-based rendezvous is used after bootstrap: once an agent has re-established its first connection (via a rendezvous agent), that reconnected friend can mediate reconnection to further friends.  Friends-based rendezvous requires at least one active connection: when all of an agent's connections break (e.g., due to an IP change), it cannot reach any friend to request mediation.  A rendezvous agent, with its known static address, provides the initial connection from which friends-based rendezvous can proceed.  Alternatively, if two agents are BLE-adjacent, friends-based rendezvous can proceed over BLE without any rendezvous agent --- BLE provides the initial connection directly.

% \dan{To add to Ohad's question: the part that should be added here is dual-stack rendezvous: friend-based rendezvous can be used over BLE if no RVs are connected; this is how the well-connected friends model work now with BLE-based signaling.}

% \udi{See corrected text above: BLE alternative now mentioned explicitly.}

%------------------------------------------------------------------------------
\subsection{Address Maintenance}
\label{sec:addr-maintenance}
%------------------------------------------------------------------------------
Each agent's friends list is extended to include IP addresses: for each friend, the agent records not only the communication channel but also the friend's current public address.  After reconnection, agents notify their friends of their new address via ordinary GLP messages.  This information is maintained as part of the social graph.


% \dan{In my opinion, GLP should be oblivious to IP addresses, and only know public keys. Just like how the IP layer does not know MAC addresses, the GLP layer should not know the addresses of the layer below. This should not be the concern of the GLP layer. If we attach the friendship attestation, the networking layer could compare public keys to know if the peers are friends.}

% \udi{Disagree.  The rendezvous agent is a GLP agent --- its coordination logic (matching reconnect/available requests, relaying addresses) is implemented in GLP code (see Appendix~\ref{appendix:rv-agent}).  Similarly, the friends-based mediation step is GLP code that looks up addresses in the friends list and sends punch instructions (Appendix~\ref{appendix:fmrv}).  Moving address handling to the networking layer would mean reimplementing this logic outside GLP, losing the advantage of having the entire rendezvous protocol be application-level GLP code.  The \code{peer\_address/2} and \code{punch\_udp/1} system predicates are the minimal interface between GLP and the networking layer; everything else is GLP.}

% \dan{I agree there is an advantage to having GLP manage friend-mediated network operations, but it cannot do all of it since it does not have any access to UDX or the network stack. I think it makes sense to have the networking layer own pure NAT traversal mechanics (send UDP probes, detect hole-punch success, manage UDX streams) while GLP owns coordination (who should punch to whom, when, with what authorization). }

% \udi{We agree with this split, and it is what the current design implements.  The system predicates are exactly this interface: \code{peer\_address/2} reads an address from the networking layer, \code{punch\_udp/1} instructs the networking layer to initiate a punch, and \code{verify\_friendship/4} delegates cryptographic verification.  GLP code decides \emph{who} should punch \emph{whom} and \emph{when}; the networking layer handles UDP probes, hole-punch detection, and stream setup.  If UDX is replaced with QUIC, only the networking-layer implementation of \code{punch\_udp} changes; the GLP code remains the same.  The predicate name could be generalized (e.g., \code{initiate\_punch/1}), but the abstraction boundary is already correct.}

% \dan{OK, then please generalize predicates: hole-punching and UDP are temporary and could be changed tomorrow. I would be happy with \code{allow\_egress} instead of sticking with hole punching, as that is intrinsic to the networking layer. GLP should not need to know about the internals of the networking layer, other than the address }

% \udi{Disagree.  This paper specifies the networking API for the current implementation, which uses UDP hole-punching through NAT.  The system predicates name the specific operations the GLP code instructs the networking layer to perform.  If the transport changes and hole-punching is no longer needed, the system predicates and the GLP code that uses them would be revised accordingly --- but abstracting away the current mechanism would obscure what the protocol actually does.}

An agent behind NAT cannot directly observe its own public address.  Instead, it learns its address from the rendezvous agent: when an agent connects to a rendezvous agent, the rendezvous agent observes the agent's public address (the source IP:port of the incoming packets) via the \code{peer\_address/2} system predicate and reports it back.  The agent then shares this address with its friends via GLP messages.  No external address-discovery infrastructure (such as STUN servers) is needed; the rendezvous agent serves this role.

%------------------------------------------------------------------------------
\subsection{Rendezvous Protocol}
\label{sec:rv-protocol}
%------------------------------------------------------------------------------
When agent~$A$'s IP address changes, all of~$A$'s connections break.  Reconnection proceeds in two phases: bootstrap via a rendezvous agent, then friends-based rendezvous for the remaining friends.

\subsubsection{Bootstrap}
\label{sec:bootstrap}

Agent~$A$ detects the IP change via \func{onConnectivityStatus} (Section~\ref{sec:discovery}).  Agent~$B$ detects that~$A$ went silent via \func{onPeerDisconnected}.  Both~$A$ and~$B$ know which rendezvous agents the other uses (learned while they were connected).
% \dan{Same comment as before: $A$ does not need to know the RV $B$ uses, but only the RV it uses to reach $B$. If accepted, change the protocol accordingly to use $A$'s set of RV agents. Since RV agents are not members of the social graph, all RV agents could gossip between themselves to find out friendship status of peers.}

% \udi{See response to the same comment in Section~\ref{sec:rv-agents}.  Both sides must converge on the same RV agent; $B$ contacts $A$'s RV agents (not $B$'s) because $A$ is the one who moved.}

\begin{enumerate}
\item $A$ connects to a rendezvous agent~$S$ (outbound to~$S$'s known static address---always succeeds through NAT).  $A$ sends a cold-call message: $\mathit{reconnect}(A_{\mathit{pk}}, B_{\mathit{pk}}, \mathit{Proof}, \mathit{Reply})$, where $\mathit{Proof}$ is~$A$'s friendship attestation for~$B$ and $\mathit{Reply}$ is a fresh reply variable.

\item $B$ detects that~$A$ disconnected.  $B$ sends a UDP packet to each of~$A$'s known rendezvous agents, opening a NAT hole so that the rendezvous agent can reach~$B$.  $B$ sends a cold-call message: $\mathit{available}(B_{\mathit{pk}}, A_{\mathit{pk}}, \mathit{Reply})$.

\item $S$ verifies~$A$'s friendship proof, observes~$A$'s public address via \code{peer\_address/2} (Section~\ref{sec:rv-predicates}), matches~$A$'s reconnect request with~$B$'s availability message, and sends~$B$ a $\mathit{punch}(\mathit{Addr}_A)$ message via~$B$'s reply variable.  If the proof is invalid, $S$ binds~$A$'s reply to $\mathit{rejected}$ and does not store the request.

\item $B$ receives $\mathit{punch}(\mathit{Addr}_A)$ and calls \code{punch\_udp/1} (Section~\ref{sec:rv-predicates}) to send UDP packets to~$A$'s new address.  Since~$A$ already has an open mapping to~$S$ (and~$B$'s NAT now has a mapping for~$A$), the packets traverse the NAT.  A direct UDP path is established.

\item The networking layer opens a UDX stream over the new UDP path and authenticates via Noise~XX (Section~\ref{sec:ip-connection}).  \func{onPeerConnected} fires at both agents.
\end{enumerate}

If the first rendezvous agent is unreachable or does not respond within a timeout,~$A$ and~$B$ proceed to the next rendezvous agent on their shared list, ordered deterministically (e.g., lexicographically by public key).  Since both sides compute the same ordered list, they converge on the same~$S$.

\subsubsection{Friends-Based Rendezvous}
\label{sec:friends-rv-protocol}

After bootstrap reconnects~$A$ to one friend~$B$, friends-based rendezvous reconnects~$A$ to the remaining friends.  When~$A$'s IP changes, all of~$A$'s connections break, but the connections among~$A$'s friends are unaffected (assuming only~$A$ changed address).  Each agent maintains a friends-of-friends map via the social graph protocol, so every friend of~$A$ knows which other agents are also friends of~$A$.

The protocol proceeds as follows:
\begin{enumerate}
\item Bootstrap reconnects~$A$ to friend~$B$.
\item $B$ knows (from its friends-of-friends map) which of its own friends are also friends of~$A$.  $B$ immediately performs friends-based mediation for every such common friend~$C$: $B$ looks up~$A$'s and~$C$'s addresses in its friends list and sends punch instructions to both sides.  No request from~$A$ is needed.
% \dan{Since $B$'s friends-of-friends map might not be up to date, I am in favor of changing the model: $A$ should explicitly request hole-punching from $B$ (or $C$, or $D$...) rather than $B$ acting proactively.}
% \udi{$B$'s friends-based mediation happens automatically for all common friends of $A$ and $B$ (step~2 above).  For friends of $A$ that are not common friends of $A$ and $B$, $A$ requests mediation from reconnected friends or goes through the RV server (step~3 above).  The friends-of-friends map need not be fully up to date --- it reflects the state at disconnection time, which is sufficient to identify common friends. See third item added below.}
% \dan{I find it hard to believe that the map instantly updates as soon as someone changes their friend list, and it needs to be: otherwise $B$ may continuously prompt $A$ and $C$ to punch for each other, despite them not being friends, and vice-versa: it may not prompt peers it should. The direction should be reversed.}

% \udi{If $B$'s map is slightly behind, the worst case is a harmless failed attempt: $C$ ignores the punch instruction, or the Noise~XX handshake fails because $A$ and $C$ are no longer friends.  $A$ also knows its own friend list, so it will not act on a \textit{punch\_to} for an agent it has unfriended.  Reversing the direction does not solve the staleness concern --- $A$'s FoF map is equally behind in the same way.  The advantage of $B$ acting proactively is reduced round-trips: $B$ mediates for all common friends in parallel immediately upon reconnection, rather than waiting for $A$ to request each one sequentially.}

\item For friends of~$A$ that are not common friends of~$A$ and~$B$, $A$ consults its friends-of-friends map.  For each such friend~$E$, if~$A$ has a reconnected friend~$D$ that is also a friend of~$E$, $A$ sends~$D$ a mediation request for~$E$.  For friends of~$A$ for which no reconnected friend can mediate, $A$ contacts the rendezvous server to reconnect to them.
\end{enumerate}

The process terminates when all of~$A$'s friends are reconnected, or when the remaining friends cannot be reached through any available mediator.

A single mediation step works as follows:
\begin{enumerate}
\item The mediating friend~$D$ looks up~$A$'s and~$C$'s current addresses in its friends list.
\item $D$ sends~$C$: $\mathit{punch\_to}(\mathit{addr}_A)$; sends~$A$: $\mathit{punch\_to}(\mathit{addr}_C)$.  If~$C$ is not in~$D$'s friends list, the $\mathit{punch\_to}$ message carries $\mathit{not\_found}$ and the recipient takes no networking action.
\item Both sides call \code{punch\_udp/1}; the UDX stream establishes.
\end{enumerate}

%------------------------------------------------------------------------------
\subsection{System Predicates}
\label{sec:rv-predicates}
%------------------------------------------------------------------------------
The rendezvous protocol requires the following system predicates, which provide the interface between GLP agent code and the networking layer:

\begin{longtable}{L{4cm} L{9.5cm}}
\toprule
\rowcolor{headerblue}
\textcolor{white}{\textbf{System Predicate}} &
\textcolor{white}{\textbf{Description}} \\
\midrule
\endhead
\code{peer\_address(PubKey, Addr)} & Unifies \code{Addr} with the public IP address and UDP port of the connected peer identified by \code{PubKey}, as observed by the networking layer from the source address of incoming packets.  Available only for peers with an active connection.  Used by rendezvous agents to learn the addresses of the agents they are helping.  Also serves as the address-discovery mechanism: an agent behind NAT learns its own public address from the rendezvous agent, which observes it via this predicate. \\
\rowcolor{rowgray}
\code{punch\_udp(Addr)} & Instructs the networking layer to begin sending UDP packets to \code{Addr} to open a NAT mapping. \\
\code{verify\_friendship(Proof, A, B, Result)} & Verifies that \code{Proof} is a valid friendship attestation proving that agents \code{A} and \code{B} are friends.  Binds \code{Result} to \code{ok} or \code{invalid}.  Used by rendezvous agents. \\
\bottomrule
\end{longtable}

All other aspects of the rendezvous protocol---the pending request table on~$S$, the coordination logic, the friends-based traversal strategy, the address maintenance in the friends list---are implemented as ordinary GLP agent code.

%------------------------------------------------------------------------------
\subsection{GLP Implementation}
\label{sec:rv-glp-impl}
%------------------------------------------------------------------------------
The rendezvous protocol has been partially implemented and verified in GLP.  We describe what has been implemented, what remains, and what is needed to complete the implementation.

\paragraph{What has been implemented.}
Three components have been implemented in GLP, pass SRSW checking and type checking, and have been verified with comprehensive test suites (presented in full in the appendices):
\begin{itemize}
\item \textbf{Rendezvous agent} (Appendix~\ref{appendix:rv-agent}): the server-side agent that processes cold-call $\mathit{reconnect}$ and $\mathit{available}$ messages, verifies friendship proofs, observes connecting agents' addresses via \code{peer\_address/2}, maintains a pending request list, and matches requesting pairs.  Invalid proofs are rejected via the reply variable.  The key mechanism is the \emph{reply-variable chain}: reply variables embedded in cold-call messages are threaded through writer$\to$reader handoffs and assigned when both sides are matched, propagating values back to the original senders.

\item \textbf{Rendezvous client} (Appendix~\ref{appendix:rv-client}): the client-side building blocks used by agents to construct cold-call messages with embedded reply variables and to dispatch on the replies ($\mathit{punch}$, $\mathit{initiated}$, or $\mathit{rejected}$).  An end-to-end test wires client and server together on a shared message stream, demonstrating the full bootstrap protocol in both arrival orders.

\item \textbf{A single friends-based mediation step} (Appendix~\ref{appendix:fmrv}): an extension of the social graph agent in which the friends list is extended to include each friend's current IP address.  The mediating friend uses a fused \code{lookup\_and\_send} procedure that, in a single traversal, locates a friend's entry, prepends a $\mathit{punch\_to}$ message to that friend's output stream, and returns the friend's address.  The $\mathit{mediate}$ clause calls this procedure twice with cross-references, leveraging GLP's suspension mechanism to resolve the mutual address dependency.  Unknown peers are handled gracefully via an $\mathit{AddrResult}$ type ($\mathit{found}(\mathit{Addr})$ or $\mathit{not\_found}$).
\end{itemize}

\paragraph{What remains.}
The friends-based rendezvous protocol (Section~\ref{sec:friends-rv-protocol}) relies on a \emph{friends-of-friends map}: each agent must know, for each friend, who that friend's friends are, in order to identify common friends for proactive mediation and to select mediators for explicit requests.  This map is not yet implemented in GLP.  The friends-of-friends protocol---in which each agent broadcasts friendship updates to all its friends via a shared broadcast stream, and each recipient maintains a local map of observed friendships---has been specified and formally verified (with safety and liveness proofs) in a companion work, but in pseudocode rather than GLP.  The current GLP social graph agent tracks only direct friends.

Without the friends-of-friends map, the implemented code covers individual components (the rendezvous agent, the client protocol, and a single mediation step) but cannot drive the full reconnection protocol described in Section~\ref{sec:friends-rv-protocol}.

\paragraph{What is needed.}
To complete the implementation, the following steps are required:
\begin{enumerate}
\item \textbf{Friends-of-friends map in GLP.}  Implement the friends-of-friends protocol as an extension of the GLP social graph agent: each agent broadcasts friendship changes to all friends and maintains a local map of each friend's friends and their status.
\item \textbf{Proactive mediation by the bootstrap friend.}  When~$B$ reconnects to~$A$ (via the rendezvous agent), $B$ uses the friends-of-friends map to identify all common friends of~$A$ and~$B$, and proactively mediates for each of them.
\item \textbf{Explicit mediation requests.}  For friends of~$A$ not covered by~$B$'s proactive mediation, $A$ uses its friends-of-friends map to identify a reconnected friend that can mediate, and sends an explicit request.
\item \textbf{Integration.}  Combine the rendezvous client, the friends-based mediation step, the friends-of-friends map, and the mediation logic into the social graph agent, so that reconnection is triggered automatically when \func{onPeerDisconnected} fires.
\end{enumerate}

%------------------------------------------------------------------------------
\subsection{Failure and Fallback}
\label{sec:rv-fallback}
%------------------------------------------------------------------------------
UDP hole-punching succeeds for cone NAT (full-cone, restricted-cone, port-restricted cone) but may fail for symmetric NAT, where the mapped port differs per destination.  When hole-punching fails, agents may fall back to relaying messages through the social graph: a connected common friend forwards madGLP messages on behalf of the unreachable peer.  Relaying is an application-level concern, implemented as GLP agent code; the networking layer does not relay.
